With the concept of hierarchical central distributions in place, we now develop the theoretical machinery needed
to characterize and compute them.
The following results provide optimization-theoretic characterizations of central distributions that form the
foundation of the subsequent analysis.

Throughout this section, all random variables are assumed to be discrete with finite support.
Before proceeding, we establish the existence of central distributions in the finite setting considered throughout
this section.

\begin{theorem}
    Let $X$ be a discrete random variable with finite support, and let \setQ be a finite set of probability
    distributions over $X$.
    Then a central distribution of \setQ for $X$ exists.
\end{theorem}

\begin{proof}
    For a fixed distribution $Q$, we may write
    \[
        \kl = c - \sum_x Q(x) \ln P(x).
    \]
    Since $\ln$ is upper semicontinuous on $[0,+\infty)$ and $X$ has finite support, it follows that
    $\kl$ is lower semicontinuous as a function of $P$.

    Because \setQ is finite, $\sFunc[P(X)][\setQ]$ is the maximum of finitely many lower semicontinuous functions
    and is therefore itself lower semicontinuous.

    The set of probability distributions over $X$ is a compact simplex.
    Hence, by the extreme value theorem, $\sFunc[][\setQ]$ attains its minimum.
    Therefore there exists a probability distribution \cDist such that
    \[
        \sFunc[\cDist(X)][\setQ] = \inf_P \sFunc[P(X)][\setQ].
    \]
    Thus $C$ is a central distribution of \setQ for $X$.
\end{proof}

The central distribution is essentially the solution to a convex--concave game.
Because solutions of convex--concave games correspond to saddle points,
by defining the appropriate objective function, we can characterize the central
distribution using standard convex optimization techniques.

\begin{definition} [Saddle Point]
    Consider a function $f:\set{w} \times \set{z} \to \field{R}$.
    We refer to a pair $(w^*, z^*)$ as a \emph{saddle point} for $f$, if we have
    \[
        f(w^*, z) \leq f(w^*, z^*) \leq f(w, z^*),
    \]
    for all \inSet{w} and \inSet{z}.
\end{definition}

Next, we establish the connection between saddle points and central distributions.
One significant insight provided by the next theorem is that it suggests inference and searching for a central
distribution can be combined and performed simultaneously when finding a saddle point.

\begin{theorem}
    % I've been unable to prove that the central distribution is a saddle point for f, or characterize conditions
    % that imply that.
    \label{th:saddle_is_central}
    Let $X$ be a discrete random variable, and \setQ be a finite set of distributions over $X$.
    Consider the following function:
    \begin{equation*}
        f(P, R)= \sum_i \kl[Q_i] R(i),
    \end{equation*}
    where $P$ and $R$ are probability distributions, $Q_i \in \setQ$ and the index $i$ ranges over all elements of \setQ.

    If \saddle is a saddle point for $f$, then \saddP is a central distribution of \setQ for $X$.
\end{theorem}

\begin{proof}
    Let \cDist be a central distribution of \setQ for $X$, and \saddle be a saddle point for $f$.
    Since $X$ is a discrete random variable and \setQ is finite, \cDist exists.
    As $C$ is a probability distribution, the point $(\cDist, \saddR)$ lies within the domain of $f$.
    By the definition of a saddle point, this implies that
    \begin{equation}
        \label{eq:th_proof_saddle_lq_dx}
        f\saddle \leq f(\cDist, \saddR).
    \end{equation}
    Because \saddR is a probability distribution, the definition of $f$ implies:
    \[
        f(\cDist, \saddR) \leq \sFunc[\cDist].
    \]

    On the other hand, by the definition of a saddle point
    \[
        f\saddle = \max_{R} f(\saddP, R),
    \]
    where $R$ is a probability distribution, and therefore
    \[
        \max_{R} f(\saddP, R) = \sFunc[\saddP].
    \]
    Thus, Equation~\ref{eq:th_proof_saddle_lq_dx} implies
    \[
        \sFunc[\saddP] \leq \sFunc[\cDist].
    \]
    Since \cDist is a central distribution, we must have \[\sFunc[\saddP] = \sFunc[\cDist],\] and $\saddP$ is also a
    central distribution of \setQ for $X$.
\end{proof}

\begin{lemma}
    \label{lm:saddle_conditions}
    Let $f$ be the function defined in Theorem~\ref{th:saddle_is_central}.
    The point \saddle is a saddle point for $f$, if and only if there are constants $\lambda$, $\eta$ and a
    function $\mu$, such that for every $i$, $x$ the following conditions hold:
    \begin{gather}
        \label{eq:kkt_lambda}
        \klExp[Q_i][\saddP][x] - \mu(i) = \lambda, \\
        \label{eq:kkt_w_avg}
        \saddP(x) = \sum_i \saddR(i) Q_i(x) \\
        \label{eq:kkt_r_sums_one}
        \sum_i \saddR(i) = 1, \\
        \label{eq:saddle_lm_r_geq_0}
        \saddR(i) \geq 0, \\
        \nonumber
        \mu(i) \geq 0, \\
        \label{eq:saddle_lm_rmu_is_zero}
        \mu(i) \saddR(i) = 0.
    \end{gather}

    Furthermore, for a fixed probability distribution $R$, a distribution $\hat{P}$ minimizes $f(\none,R)$, if and
    only if for all $x$ it satisfies
    \begin{equation}
        \label{eq:saddle_lm_minimizer_dist}
        \hat{P}(x) = \sum_i Q_i(x)R(i).
    \end{equation}
\end{lemma}

\begin{proof}
    Similar to $f$, we define $\hat{f}$ as follows:
    \begin{equation*}
        \hat{f}(P, R) = \fExp,
    \end{equation*}
    where the constraints that $P$ and $R$ be probability distributions are omitted.
    This enables us to formulate the saddle point of $f$ as the solution to a constrained optimization problem
    defined for $\hat{f}$.

    Let \saddle be a simultaneous solution to the following two optimization problems:

    \begin{equation}
        \begin{aligned}
            &\text{minimize} & &\hat{f}(P,R) & &\text{w.r.t.} \ P \\
            &\text{subject to} & &\sum_{x} P(x) = 1, \\
            \\
            &\text{maximize} & &\hat{f}(P,R) & &\text{w.r.t.} \ R \\
            &\text{subject to} & &R(i) \geq 0, & &\text{for all} \ i  \\
            & & &\sum_{i} R(i) = 1 &
        \end{aligned}\label{eq:twin_optimization}
    \end{equation}
    Clearly, if $\saddle$ exists, it will be a saddle point for $f$.

    Both of these optimization problems are convex and differentiable in the sense of convex
    optimization; see~\cite{boyd2004convex}.
    To establish this, consider that all the constraint functions are affine, and therefore are convex, concave and
    differentiable.

    Additionally, $\hat{f}$ is a differentiable function with respect to both $R$ and $P$.
    For any fixed $P$, $\hat{f}(P,\none)$ is an affine function and is concave.

    On the other hand, for a fixed solution to the second optimization problem denoted by $\hat{R}$, we can write
    \begin{equation}
        \label{eq:simplified_f}
        \hat{f}(P,\hat{R}) = c - \sum_{i,x} \hat{R}(i) Q_i(x) \ln P(x),
    \end{equation}
    where $c$ is not a function of $P$.
    Since $\hat{R}$ is a solution to the second optimization problem, it satisfies the corresponding constraints and
    is therefore a valid probability distribution.
    Hence, for all $i$, $x$ we have $\hat{R}(i)Q_i(x) \geq 0$, and $\hat{f}(\none, \hat{R})$  is
    the negative of a linear combination of concave $\ln P(x)$ functions with nonnegative coefficients.
    This function is convex.

    Thus, both optimization problems are convex and differentiable, with affine constraints.
    The constraints simply restrict the feasible set to the set of probability distributions, and it is always possible
    to find feasible points that meet all the constraints.
    Therefore, Slater’s condition is satisfied, and the KKT conditions provide necessary
    and sufficient conditions for optimality~\cite{boyd2004convex}.
    In other words, if there exists a point $\saddle$ that satisfies the KKT conditions for both problems, it will be a
    saddle point for $f$, and vice versa.

    The KKT conditions are as follows for every $x$, $i$ and constants $\lambda$, $\eta$:
    \begin{gather}
        \label{eq:kkt_proof_R_is_pxIe}
        \sum_i \saddR(i) Q_i(x) = \eta \saddP(x) \\
        \label{eq:kkt_proof_p_sums_one}
        \sum_{x} \saddP(x) = 1, \\
        \nonumber
        \klExp[Q_i][\saddP][x] - \mu(i) = \lambda, \\
        \label{eq:kkt_proof_r_sums_one}
        \sum_i \saddR(i) = 1, \\
        \nonumber
        \saddR(i) \geq 0, \\
        \nonumber
        \mu(i) \geq 0, \\
        \nonumber
        \mu(i) \saddR(i) = 0.
    \end{gather}

    By summing Equation~\ref{eq:kkt_proof_R_is_pxIe} over $x$ and using Equation~\ref{eq:kkt_proof_r_sums_one},
    we have
    \[
        \eta = 1\,.
    \]
    Thus, we get Equation~\ref{eq:kkt_w_avg} which also implicitly implies Equation~\ref{eq:kkt_proof_p_sums_one}.

    To prove the second part of the lemma,
    we observe that when $R$ is a probability distribution, Equations~\ref{eq:kkt_proof_R_is_pxIe}
    and~\ref{eq:kkt_proof_p_sums_one} are the KKT conditions for the minimization problem
    in Equation~\ref{eq:twin_optimization}.
    Since this minimization problem is convex and satisfies Slater's condition,
    the KKT conditions are necessary and sufficient.
    Hence
    \[
        \hat P(x)=\sum_i Q_i(x)R(i)
    \]
    if and only if $\hat P$ minimizes $f(\none,R)$, where $R$ is any fixed distribution.
\end{proof}

\begin{lemma}
    \label{lm:saddle-is-maximizer}
    Let \setQ be a finite set of probability distributions over a discrete random variable $X$.
    Define the function $g$ as:
    \[
        g(R) = \sum_i \kl[Q_i][R] R(i),
    \]
    where $Q_i \in \setQ$, and $R$ is a probability distribution such that \[R(X, I) = Q_i(X)R(I).\]
    Let $f$ be the function defined in Theorem~\ref{th:saddle_is_central}.
    If \saddle is a saddle point for $f$, then the joint distribution $Q_i(X)\saddR(I)$ is a maximizer of $g$.
\end{lemma}

\begin{proof}
    Let \saddle be a saddle point for $f$.
    By Lemma~\ref{lm:saddle_conditions}, \saddP satisfies Equation~\ref{eq:kkt_w_avg}.
    Letting $\saddR(X, I) = Q_i(X) \saddR(I)$ we have
    \[
        g(\saddR) = f\saddle.
    \]
    Now, let $R$ be an arbitrary probability distribution such that $g(R)$ is defined.
    From the second part of Lemma~\ref{lm:saddle_conditions}, it follows that $g(R)$ indeed corresponds to the
    minimum of $f(\none, R)$, which implies
    \[
        g(R) \leq f(\saddP,R).
    \]

    On the other hand, since \saddle is a saddle point, we also have
    \[
        f(\saddP, R) \leq f\saddle.
    \]
    Combining the two inequalities, we get:
    \[
        g(R) \leq g(\saddR).
    \]
    This shows that $g(\saddR)$ is the maximum of $g$, as desired.
\end{proof}

The following theorem gives a standard sufficient condition for optimality in the KL-minimax setting which is
the rationale behind the term ``central distribution.''
Like the center of a circle, the central distribution is equidistant from the boundary.

\begin{theorem} [Centrality]
    \label{th:centrality}
    Let $X$ be a discrete random variable, and let \setQ be a finite set of distributions over $X$.
    Suppose \cDist is a convex combination of elements of \setQ, with all mixture weights strictly positive.
    If there exists a constant $\lambda$ such that
    \begin{equation}
        \label{eq:equal-div}
        \kl[][\cDist] = \lambda,
    \end{equation}
    for every \inSet{Q}, then \cDist is the unique central distribution of \setQ for $X$.
\end{theorem}

\begin{proof}
    Lemma~\ref{lm:saddle_conditions} proves a set of sufficient conditions for a saddle point.
    Note that more restrictive conditions can also be imposed to derive an alternative set of sufficient conditions.
    Suppose additionally that $\saddR(i) > 0$ for every $i$.
    Then Equation~\ref{eq:saddle_lm_rmu_is_zero} implies $\mu(i) = 0$,
    and Lemma~\ref{lm:saddle_conditions} conditions can be simplified to
    \begin{gather*}
        \klExp[Q][\saddP][x] = \lambda, \\
        \saddP(x) = \sum_i \saddR(i) Q_i(x),
    \end{gather*}
    which must hold for every \inSet{Q} and some strictly positive probability distribution \saddR.
    Since \cDist is a convex combination of elements of \setQ, the corresponding convex
    weights define the distribution \saddR and \cDist satisfies these conditions.

    Therefore, by Lemma~\ref{lm:saddle_conditions} and Theorem~\ref{th:saddle_is_central}, \cDist is a central
    distribution of \setQ for $X$.

    On the other hand, for any probability distribution $C'$, we have
    \begin{align}\label{eq:cd-relation-to-lambda}
        \nonumber
        \kl[\cDist][C']
        &= \sum_i \saddR(i)\left(\kl[Q_i][C'] - \kl[Q_i][\cDist] \right) \\
        \nonumber
        &= \sum_i \saddR(i)\kl[Q_i][C'] - \lambda \\
        &\leq \sFunc[C'] - \lambda.
    \end{align}
    If $C'$ is a central distribution of \setQ for $X$,
    the definition of a central distribution implies
    \[
        \sFunc[C'] = \sFunc[\cDist] = \lambda.
    \]
    Therefore, by non-negativity of the KL divergence, we have
    \[
        \kl[\cDist][C'] = 0.
    \]
    Hence \cDist and $C'$ are identical probability distributions.
\end{proof}

\begin{corollary}
    \label{cor:C-maximizes-g}
    Under the assumptions of Theorem~\ref{th:centrality}, \cDist is a maximizer of the function $g$ defined in
    Lemma~\ref{lm:saddle-is-maximizer}.
\end{corollary}

\begin{proof}
    In the proof of Theorem~\ref{th:centrality}, we showed that $(\cDist, \saddR)$ is a saddle point for $f$.
    By letting $C(X, I) = Q_i(X)\saddR(I)$, Lemma~\ref{lm:saddle-is-maximizer} implies that \cDist is a maximizer
    of the function $g$.
\end{proof}

\begin{lemma}
    \label{lm:imperfection}
    Assuming the conditions of Theorem~\ref{th:centrality} and using the same $\lambda$,
    if $P$ is a convex combination of elements of \setQ, then there exists some $Q_0 \in \setQ$ such that
    \begin{equation*}
        \kl[Q_0][] \leq \lambda.
    \end{equation*}
\end{lemma}

\begin{proof}
    We proceed by contradiction.
    Suppose that for every \inSet{Q}, we have $\kl > \lambda$.
    This implies that
    \[
        g(P) > \lambda = g(C),
    \]
    where $C$ is the distribution specified in Theorem~\ref{th:centrality}, and $g$ is the function defined in
    Lemma~\ref{lm:saddle-is-maximizer}.
    Note that $g(P)$ is defined, because $P$ is a convex combination of elements of \setQ.

    However, by Corollary~\ref{cor:C-maximizes-g}, the distribution $C$ maximizes $g$, which contradicts the
    inequality $g(P) > g(C)$.
\end{proof}
